\section{Randomized --- Space Amplify Poly Machine}
%%% generated by the blueprint pipeline from TCSlib.Complexity.Randomized.SpaceAmplifyPolyMachine
\subsection*{Randomized --- Space Amplify Poly Machine: Introduction}

Declarations of the Lean module \texttt{TCSlib.Complexity.Randomized.SpaceAmplifyPolyMachine}.

\subsection*{Randomized --- Space Amplify Poly Machine: Results}

\begin{definition}[Register copy sources]
\lean{Randomized.PolyAmp.CpS}
\label{Randomized.PolyAmp.CpS}
\leanok
Let the finite set of eight source labels consist of $nt$, $la1$, $la2$, $kR$, $ri$,
$kj$, $kc$, and $cc$.
\end{definition}

\begin{definition}[Source register assignment]
\lean{Randomized.PolyAmp.cpSrc}
\label{Randomized.PolyAmp.cpSrc}
\leanok
\uses{Randomized.PolyAmp.CpS}
Define a map from the eight source labels $nt$, $la1$, $la2$, $kR$, $ri$, $kj$, $kc$,
and $cc$ to the registers numbered $0$ through $14$ by sending them, respectively, to
registers $0$, $1$, $1$, $4$, $5$, $4$, $4$, and $9$.
\end{definition}

\begin{definition}[Destination register assignments]
\lean{Randomized.PolyAmp.cpDst}
\label{Randomized.PolyAmp.cpDst}
\leanok
\uses{Randomized.PolyAmp.CpS}
Define a map from the eight symbols
$	exttt{nt},	exttt{la1},	exttt{la2},	exttt{kR},	exttt{ri},	exttt{kj},	exttt{kc},	exttt{cc}$
to the registers numbered $0$ through $14$ by sending them, respectively, to registers
$12,13,13,5,6,10,9,8$.
\end{definition}

\begin{definition}[Machine labels]
\lean{Randomized.PolyAmp.Lab}
\label{Randomized.PolyAmp.Lab}
\leanok
\uses{Randomized.PolyAmp.CpS}
For natural numbers $cS$ and $F$, define the label type with labels
$	exttt{len},	exttt{lenI}$;
$	exttt{cp0}(k),	exttt{cp1}(k),	exttt{cp2}(k),	exttt{cp3}(k),	exttt{cp4}(k),	exttt{cp5}(k),	exttt{cp6}(k)$
for each of the eight symbols $k$; $	exttt{hT},	exttt{hH},	exttt{hB},	exttt{lI}$;
$	exttt{sT0},	exttt{sD0},	exttt{sI0}(i)$ for $i<cS$;
$	exttt{sT1},	exttt{sD1},	exttt{sI1}(i)$ for $i<F$; and the remaining labels
$	exttt{kI},	exttt{sh0}$ through
$	exttt{sh8},	exttt{rI},	exttt{rT},	exttt{rD},	exttt{iC},	exttt{iT},	exttt{iI},	exttt{sT},	exttt{sD},	exttt{coin},	exttt{up},	exttt{uT},	exttt{uI},	exttt{dn},	exttt{dD},	exttt{dT},	exttt{mv},	exttt{chk},	exttt{cI},	exttt{fin},	exttt{ad1},	exttt{ad2},	exttt{ad3},	exttt{fC},	exttt{lt1},	exttt{lt2},	exttt{ltb},	exttt{acc},	exttt{rej}$.
\end{definition}

\begin{definition}[Label assigned to each symbol]
\lean{Randomized.PolyAmp.cpEx}
\label{Randomized.PolyAmp.cpEx}
\leanok
\uses{Randomized.PolyAmp.CpS, Randomized.PolyAmp.Lab}
For natural numbers $cS$ and $F$, define a map from the eight symbols to labels by
sending the symbols to their corresponding labels as follows:
$\mathrm{nt}\mapsto\mathrm{hT}$, $\mathrm{la1}\mapsto\mathrm{sT0}$,
$\mathrm{la2}\mapsto\mathrm{sT1}$, $\mathrm{kR}\mapsto\mathrm{rI}$,
$\mathrm{ri}\mapsto\mathrm{rT}$, $\mathrm{kj}\mapsto\mathrm{sT}$,
$\mathrm{kc}\mapsto\mathrm{dD}$, and $\mathrm{cc}\mapsto\mathrm{sT}$.
\end{definition}

\begin{definition}[First indexed label with fallback]
\lean{Randomized.PolyAmp.sL0}
\label{Randomized.PolyAmp.sL0}
\leanok
\uses{Randomized.PolyAmp.Lab}
For natural numbers $cS,F$, and any natural number $i$, define a label to be
$	exttt{sI0}(i)$ when $i<cS$, and $	exttt{sT0}$ otherwise.
\end{definition}

\begin{definition}[Second indexed label with fallback]
\lean{Randomized.PolyAmp.sL1}
\label{Randomized.PolyAmp.sL1}
\leanok
\uses{Randomized.PolyAmp.Lab}
For natural numbers $cS,F$, and any natural number $i$, define a label to be
$	exttt{sI1}(i)$ when $i<F$, and $	exttt{sT1}$ otherwise.
\end{definition}

\begin{definition}[Amplification program transition function]
\lean{Randomized.PolyAmp.ampA}
\label{Randomized.PolyAmp.ampA}
\leanok
\uses{Complexity.LogProg.NARM, Randomized.PolyAmp.Lab, Randomized.PolyAmp.cpDst, Randomized.PolyAmp.cpEx, Randomized.PolyAmp.cpSrc, Randomized.PolyAmp.sL0, Randomized.PolyAmp.sL1}
For natural numbers $cS$ and $F$ and a Boolean $os$, the program with 15 registers, four
deciders, and labels parameterized by $cS$ and $F$ is defined by the transitions in the
formal statement: its labels specify deterministic register operations, branches, and
calls, as well as a choice between two destinations at the sampling label; at the
terminal labels it returns true or false, with the transition at the final decision
label determined by $os$.
\end{definition}

\begin{lemma}[Amplification Program Preconditions]
\lean{Randomized.PolyAmp.ampA_npreS}
\label{Randomized.PolyAmp.ampA_npreS}
\leanok
\uses{Complexity.LogProg.AConf, Complexity.LogProg.NARM.detPart, Complexity.LogProg.NPreS, Complexity.LogProg.PreS, Randomized.PolyAmp.Lab, Randomized.PolyAmp.ampA}
\difficulty{4}
For all natural numbers $cS,F$, Boolean values $os$, Boolean input words $x$, and
abstract configurations $a$ with fifteen registers and labels indexed by $cS$ and $F$,
the amplification program with each nondeterministic choice resolved in favor of false
satisfies its precondition at $a$ on input $x$.
\end{lemma}

\begin{definition}[Machine language tuple]
\lean{Randomized.PolyAmp.ampAs}
\label{Randomized.PolyAmp.ampAs}
\leanok
\uses{Complexity.Codec.acceptLang, Complexity.Codec.initLang, Complexity.Codec.stepLang, Complexity.inputLenLang, Turing.FinNDTM}
Given a finite nondeterministic Turing machine $N$ over the Boolean alphabet, define a
function on the four indices $0,1,2,3$ whose values, in order, are the language of valid
input lengths, the language of initial configurations of $N$, the language of valid
single-step transitions of $N$, and the language of accepting configurations of $N$.
\end{definition}

\begin{theorem}[Amplification languages are logarithmic-space decidable]
\lean{Randomized.PolyAmp.ampAs_mem}
\label{Randomized.PolyAmp.ampAs_mem}
\leanok
\uses{Complexity.Codec.acceptLang_mem, Complexity.Codec.initLang_mem, Complexity.Codec.stepLang_mem, Complexity.LOGSPACE, Complexity.inputLenLang_mem, Randomized.PolyAmp.ampAs, Turing.FinNDTM}
\difficulty{4}
For every finite-state nondeterministic Turing machine $N$ over the Boolean alphabet and
every index $j\in\{0,1,2,3\}$, the $j$th language in the four-language family associated
with $N$ belongs to logarithmic space.
\end{theorem}

\begin{definition}[Language membership indicator]
\lean{Randomized.PolyAmp.ampOr}
\label{Randomized.PolyAmp.ampOr}
\leanok
\uses{Randomized.PolyAmp.ampAs, Turing.FinNDTM, Turing.MultiTapeTM.indicator}
Given a finite nondeterministic Turing machine $N$ over the Boolean alphabet, define a
Boolean-valued function of an index $j\in\mathrm{Fin}(4)$ and a Boolean word $V$ that is true
exactly when $V$ belongs to the language at index $j$ in the four-language tuple
associated with $N$.
\end{definition}

\begin{definition}[Register state]
\lean{Randomized.PolyAmp.Reg}
\label{Randomized.PolyAmp.Reg}
\leanok
A register state is a tuple of fifteen natural numbers: the input length, the parameter
$\ell$, the radius, the exponent, $K$, the number of runs, the runs remaining, the
number of accepting runs, the current code, a candidate code, the coins remaining, and
four scratch values (for copies and a power, digit counting and comparison, scalings,
and twice the count, respectively).
\end{definition}

\begin{definition}[Register vector]
\lean{Randomized.PolyAmp.Reg.v}
\label{Randomized.PolyAmp.Reg.v}
\leanok
\uses{Randomized.PolyAmp.Reg}
For a register state, its register vector is the length-fifteen vector of natural
numbers whose entries, in order, are the input length, $\ell$, the radius, the exponent,
$K$, the number of runs, the runs remaining, the accepting-run count, the current code,
the candidate code, the coins remaining, and the four scratch values.
\end{definition}

\begin{definition}[Bounded register values]
\lean{Randomized.PolyAmp.Reg.Le}
\label{Randomized.PolyAmp.Reg.Le}
\leanok
\uses{Randomized.PolyAmp.Reg}
A register state is bounded by a natural number $B$ when each of its fifteen register
values is at most $B$.
\end{definition}

\begin{lemma}[Uniform bound on register values]
\lean{Randomized.PolyAmp.Reg.le_v}
\label{Randomized.PolyAmp.Reg.le_v}
\leanok
\uses{Randomized.PolyAmp.Reg, Randomized.PolyAmp.Reg.Le, Randomized.PolyAmp.Reg.v}
\difficulty{4}
Let $\sigma$ be a register state and $B$ a natural number. If every component of
$\sigma$ is at most $B$, then every entry of its length-fifteen register vector is at
most $B$.
\end{lemma}

\begin{lemma}[First Register Value]
\lean{Randomized.PolyAmp.Reg.v_0}
\label{Randomized.PolyAmp.Reg.v_0}
\leanok
\uses{Randomized.PolyAmp.Reg.v}
\difficulty{1}
For any register state $\sigma$ and any natural number $y$, the value of its register
vector at coordinate $0$ is its input-length field.
\end{lemma}

\begin{lemma}[Second Register Value]
\lean{Randomized.PolyAmp.Reg.v_1}
\label{Randomized.PolyAmp.Reg.v_1}
\leanok
\uses{Randomized.PolyAmp.Reg.v}
\difficulty{1}
For any register state $\sigma$ and any natural number $y$, the value of its register
vector at coordinate $1$ is its $\ell$ field.
\end{lemma}

\begin{lemma}[Third Register Value]
\lean{Randomized.PolyAmp.Reg.v_2}
\label{Randomized.PolyAmp.Reg.v_2}
\leanok
\uses{Randomized.PolyAmp.Reg.v}
\difficulty{1}
For every register state $\sigma$ and natural number $y$, the register vector of
$\sigma$ has value $s$ at coordinate $2$.
\end{lemma}

\begin{lemma}[Fourth Register Value]
\lean{Randomized.PolyAmp.Reg.v_3}
\label{Randomized.PolyAmp.Reg.v_3}
\leanok
\uses{Randomized.PolyAmp.Reg.v}
\difficulty{1}
For every register state $\sigma$ and natural number $y$, the register vector of
$\sigma$ has value $e$ at coordinate $3$.
\end{lemma}

\begin{lemma}[Register Vector Entry Four]
\lean{Randomized.PolyAmp.Reg.v_4}
\label{Randomized.PolyAmp.Reg.v_4}
\leanok
\uses{Randomized.PolyAmp.Reg.v}
\difficulty{1}
For every register state $\sigma$, the entry at position $4$ of its register vector is
its $K$ register.
\end{lemma}

\begin{lemma}[Register Vector Entry Five]
\lean{Randomized.PolyAmp.Reg.v_5}
\label{Randomized.PolyAmp.Reg.v_5}
\leanok
\uses{Randomized.PolyAmp.Reg.v}
\difficulty{1}
For every register state $\sigma$, the entry at position $5$ of its register vector is
its run-count register $R$.
\end{lemma}

\begin{lemma}[Register Vector Entry Six]
\lean{Randomized.PolyAmp.Reg.v_6}
\label{Randomized.PolyAmp.Reg.v_6}
\leanok
\uses{Randomized.PolyAmp.Reg.v}
\difficulty{1}
For every register state $\sigma$, the entry at position $6$ of its register vector is
its remaining-runs register $i$.
\end{lemma}

\begin{lemma}[Register Vector Entry Seven]
\lean{Randomized.PolyAmp.Reg.v_7}
\label{Randomized.PolyAmp.Reg.v_7}
\leanok
\uses{Randomized.PolyAmp.Reg.v}
\difficulty{1}
For every register state $\sigma$, the entry at position $7$ of its register vector is
its accepting-run count register $\mathrm{cnt}$.
\end{lemma}

\begin{lemma}[Register Vector Entry Eight]
\lean{Randomized.PolyAmp.Reg.v_8}
\label{Randomized.PolyAmp.Reg.v_8}
\leanok
\uses{Randomized.PolyAmp.Reg.v}
\difficulty{1}
For every register state $\sigma$, the entry at position $8$ of its register vector is
its current-code register $c$.
\end{lemma}

\begin{lemma}[Register Vector Entry Nine]
\lean{Randomized.PolyAmp.Reg.v_9}
\label{Randomized.PolyAmp.Reg.v_9}
\leanok
\uses{Randomized.PolyAmp.Reg.v}
\difficulty{1}
For every register state $\sigma$, the entry at position $9$ of its register vector is
its candidate-code register $d$.
\end{lemma}

\begin{lemma}[Eleventh Register Value]
\lean{Randomized.PolyAmp.Reg.v_10}
\label{Randomized.PolyAmp.Reg.v_10}
\leanok
\uses{Randomized.PolyAmp.Reg.v}
\difficulty{1}
For every register state $\sigma$ and natural number $y$, the register vector of
$\sigma$ has value $j$ at coordinate $10$.
\end{lemma}

\begin{lemma}[Twelfth Register Value]
\lean{Randomized.PolyAmp.Reg.v_11}
\label{Randomized.PolyAmp.Reg.v_11}
\leanok
\uses{Randomized.PolyAmp.Reg.v}
\difficulty{1}
For every register state $\sigma$ and natural number $y$, the register vector of
$\sigma$ has value $z$ at coordinate $11$.
\end{lemma}

\begin{lemma}[Thirteenth Register Value]
\lean{Randomized.PolyAmp.Reg.v_12}
\label{Randomized.PolyAmp.Reg.v_12}
\leanok
\uses{Randomized.PolyAmp.Reg.v}
\difficulty{1}
For every register state $\sigma$ and natural number $y$, the register vector of
$\sigma$ has value $t$ at coordinate $12$.
\end{lemma}

\begin{lemma}[Fourteenth Register Value]
\lean{Randomized.PolyAmp.Reg.v_13}
\label{Randomized.PolyAmp.Reg.v_13}
\leanok
\uses{Randomized.PolyAmp.Reg.v}
\difficulty{1}
For every register state $\sigma$ and natural number $y$, the register vector of
$\sigma$ has value $a$ at coordinate $13$.
\end{lemma}

\begin{lemma}[Fifteenth Register Value]
\lean{Randomized.PolyAmp.Reg.v_14}
\label{Randomized.PolyAmp.Reg.v_14}
\leanok
\uses{Randomized.PolyAmp.Reg.v}
\difficulty{1}
For every register state $\sigma$ and natural number $y$, the register vector of
$\sigma$ has value $w$ at coordinate $14$.
\end{lemma}

\begin{lemma}[Updating the input-length register]
\lean{Randomized.PolyAmp.Reg.u_0}
\label{Randomized.PolyAmp.Reg.u_0}
\leanok
\uses{Randomized.PolyAmp.Reg.v}
\difficulty{4}
For a register state $\sigma$ and a natural number $y$, updating entry $0$ of its
register vector to $y$ gives the register vector of the state obtained by replacing the
input-length component of $\sigma$ by $y$.
\end{lemma}

\begin{lemma}[Updating the second register]
\lean{Randomized.PolyAmp.Reg.u_1}
\label{Randomized.PolyAmp.Reg.u_1}
\leanok
\uses{Randomized.PolyAmp.Reg.v}
\difficulty{4}
For a register state $\sigma$ and a natural number $y$, updating entry $1$ of its
register vector to $y$ gives the register vector of the state obtained by replacing the
parameter $\ell$ in $\sigma$ by $y$.
\end{lemma}

\begin{lemma}[Updating the radius register]
\lean{Randomized.PolyAmp.Reg.u_2}
\label{Randomized.PolyAmp.Reg.u_2}
\leanok
\uses{Randomized.PolyAmp.Reg.v}
\difficulty{4}
For a register state $\sigma$ and a natural number $y$, replacing coordinate $2$ of its
register vector by $y$ gives the register vector of the state obtained by replacing its
radius by $y$.
\end{lemma}

\begin{lemma}[Updating the exponent register]
\lean{Randomized.PolyAmp.Reg.u_3}
\label{Randomized.PolyAmp.Reg.u_3}
\leanok
\uses{Randomized.PolyAmp.Reg.v}
\difficulty{4}
For a register state $\sigma$ and a natural number $y$, replacing coordinate $3$ of its
register vector by $y$ gives the register vector of the state obtained by replacing its
exponent by $y$.
\end{lemma}

\begin{lemma}[Updating the $K$ register]
\lean{Randomized.PolyAmp.Reg.u_4}
\label{Randomized.PolyAmp.Reg.u_4}
\leanok
\uses{Randomized.PolyAmp.Reg.v}
\difficulty{4}
For a register state $\sigma$ and a natural number $y$, replacing coordinate $4$ of its
register vector by $y$ gives the register vector of the state obtained by replacing its
$K$ value by $y$.
\end{lemma}

\begin{lemma}[Updating the Number-of-Runs Register]
\lean{Randomized.PolyAmp.Reg.u_5}
\label{Randomized.PolyAmp.Reg.u_5}
\leanok
\uses{Randomized.PolyAmp.Reg.v}
\difficulty{4}
For a register state $\sigma$ and a natural number $y$, replacing coordinate $5$ of its
register vector by $y$ gives the register vector of the state obtained by setting its
number-of-runs field to $y$.
\end{lemma}

\begin{lemma}[Updating the Runs-Remaining Register]
\lean{Randomized.PolyAmp.Reg.u_6}
\label{Randomized.PolyAmp.Reg.u_6}
\leanok
\uses{Randomized.PolyAmp.Reg.v}
\difficulty{4}
For a register state $\sigma$ and a natural number $y$, replacing coordinate $6$ of its
register vector by $y$ gives the register vector of the state obtained by setting its
runs-remaining field to $y$.
\end{lemma}

\begin{lemma}[Updating the Accepting-Runs Register]
\lean{Randomized.PolyAmp.Reg.u_7}
\label{Randomized.PolyAmp.Reg.u_7}
\leanok
\uses{Randomized.PolyAmp.Reg.v}
\difficulty{4}
For a register state $\sigma$ and a natural number $y$, replacing coordinate $7$ of its
register vector by $y$ gives the register vector of the state obtained by setting its
accepting-runs field to $y$.
\end{lemma}

\begin{lemma}[Updating the Current-Code Register]
\lean{Randomized.PolyAmp.Reg.u_8}
\label{Randomized.PolyAmp.Reg.u_8}
\leanok
\uses{Randomized.PolyAmp.Reg.v}
\difficulty{4}
For a register state $\sigma$ and a natural number $y$, replacing coordinate $8$ of its
register vector by $y$ gives the register vector of the state obtained by setting its
current-code field to $y$.
\end{lemma}

\begin{lemma}[Updating the Candidate-Code Register]
\lean{Randomized.PolyAmp.Reg.u_9}
\label{Randomized.PolyAmp.Reg.u_9}
\leanok
\uses{Randomized.PolyAmp.Reg.v}
\difficulty{4}
For a register state $\sigma$ and a natural number $y$, replacing coordinate $9$ of its
register vector by $y$ gives the register vector of the state obtained by setting its
candidate-code field to $y$.
\end{lemma}

\begin{lemma}[Updating the coins register]
\lean{Randomized.PolyAmp.Reg.u_10}
\label{Randomized.PolyAmp.Reg.u_10}
\leanok
\uses{Randomized.PolyAmp.Reg.v}
\difficulty{4}
For a register state $\sigma$ and a natural number $y$, updating entry $10$ of its
register vector to $y$ gives the register vector of the state obtained by replacing the
number of coins remaining in $\sigma$ by $y$.
\end{lemma}

\begin{lemma}[Updating the copies and power register]
\lean{Randomized.PolyAmp.Reg.u_11}
\label{Randomized.PolyAmp.Reg.u_11}
\leanok
\uses{Randomized.PolyAmp.Reg.v}
\difficulty{4}
For a register state $\sigma$ and a natural number $y$, replacing coordinate $11$ of its
register vector by $y$ gives the register vector of the state obtained by replacing its
copies-and-power scratch value by $y$.
\end{lemma}

\begin{lemma}[Updating the digit-count register]
\lean{Randomized.PolyAmp.Reg.u_12}
\label{Randomized.PolyAmp.Reg.u_12}
\leanok
\uses{Randomized.PolyAmp.Reg.v}
\difficulty{4}
For a register state $\sigma$ and a natural number $y$, replacing coordinate $12$ of its
register vector by $y$ gives the register vector of the state obtained by replacing its
digit-count-and-comparison scratch value by $y$.
\end{lemma}

\begin{lemma}[Updating the scaling register]
\lean{Randomized.PolyAmp.Reg.u_13}
\label{Randomized.PolyAmp.Reg.u_13}
\leanok
\uses{Randomized.PolyAmp.Reg.v}
\difficulty{4}
For a register state $\sigma$ and a natural number $y$, replacing coordinate $13$ of its
register vector by $y$ gives the register vector of the state obtained by replacing its
scaling scratch value by $y$.
\end{lemma}

\begin{lemma}[Updating the twice-count register]
\lean{Randomized.PolyAmp.Reg.u_14}
\label{Randomized.PolyAmp.Reg.u_14}
\leanok
\uses{Randomized.PolyAmp.Reg.v}
\difficulty{4}
For a register state $\sigma$ and a natural number $y$, replacing coordinate $14$ of its
register vector by $y$ gives the register vector of the state obtained by replacing its
twice-count value by $y$.
\end{lemma}

\begin{definition}[Amplified-machine configuration probability]
\lean{Randomized.PolyAmp.PS}
\label{Randomized.PolyAmp.PS}
\leanok
\uses{Complexity.LogProg.PSem, Complexity.LogProg.RegLe, Randomized.PolyAmp.Lab, Randomized.PolyAmp.Reg, Randomized.PolyAmp.Reg.v, Randomized.PolyAmp.ampA, Randomized.PolyAmp.ampOr, Turing.FinNDTM}
Let $N$ be a finite nondeterministic Turing machine over the Boolean alphabet, let $os$
be a Boolean, let $x$ be a finite Boolean input, let $B$ be a natural number, let $l$ be
a label, let $\sigma$ be a register state, and let $p$ be rational. The predicate holds
when the probability of the configuration with label $l$, register vector of $\sigma$,
and no output, on input $x$, in the machine obtained by applying the amplification
construction with parameters $cS,F,os$ to the disjunction of $N$, satisfies the register
bound $B$, is $p$.
\end{definition}

\begin{lemma}[Invariance Under Register-State Equality]
\lean{Randomized.PolyAmp.PS.of_eq}
\label{Randomized.PolyAmp.PS.of_eq}
\leanok
\uses{Randomized.PolyAmp.Lab, Randomized.PolyAmp.PS, Randomized.PolyAmp.Reg, Turing.FinNDTM}
\difficulty{1}
Fix a finite nondeterministic Turing machine $N$ with Boolean input symbols, a Boolean
parameter $o$, a Boolean input word $x$, a natural-number register bound $B$, a control
label $l$, and a rational number $p$. For register states $\sigma$ and $\sigma'$, if
$\sigma=\sigma'$ and the polynomial-amplification program initialized at $l$ with
register state $\sigma'$ has all registers bounded by $B$, terminates for every sequence
of choice bits, and has probabilities $p$ and $1-p$ of returning true and false,
respectively, at every choice-sequence length for which all branches halt, then the same
properties hold when it is initialized with $\sigma$.
\end{lemma}
